\chapter{ESTADO DEL ARTE}
\label{chap:estado_arte}

A lo largo del tiempo, la evaluación del biodeterioro en monumentos arqueológicos e históricos ha pasado por distintas etapas tecnológicas, donde cada método presentó limitaciones operativas y alcances técnicos que influyeron en la evolución del área. En este contexto, el propósito de este capítulo es revisar de manera crítica las soluciones analíticas y de intervención desarrolladas tanto en el ámbito académico como industrial, estableciendo un marco comparativo basado en parámetros de rendimiento, cobertura y viabilidad técnica. Esto permitirá comprender las fortalezas y limitaciones de los enfoques previos, así como justificar el desarrollo y la adopción de aproximaciones diagnósticas contemporáneas. 

\section{Soluciones precedentes}
\label{sec:soluciones_precedentes}
\subsection{Protocolos basados en cultivo microbiológico (método tradicional)}
\label{subsec:cultivo_tradicional}
Durante décadas, el estudio de los ecosistemas microbianos en el patrimonio pétreo dependió casi exclusivamente de técnicas dependientes de cultivo. Las muestras extraídas del sustrato se sembraban en medios de laboratorio generales, como Agar Nutritivo, MEA o PDA, frecuentemente suplementados con antibióticos o cicloheximida para la selección de grupos específicos \cite{branysova2022}. Este enfoque permitió aislar e identificar morfológicamente a numerosos taxa bacterianos y fúngicos asociados al biodeterioro.

No obstante, esta metodología introduce sesgos selectivos, siendo su principal limitación la estructural: se estima que aproximadamente el 99 \% de los microorganismos ambientales son "viables pero no cultivables" (VBNC) en condiciones de laboratorio estándar, dado que los medios artificiales no replican las condiciones oligotróficas, de estrés hídrico y radiación UV características de las superficies pétreas \cite{branysova2022, sterflinger2013, suchy2024}. En contraste con las tecnologías moleculares modernas, los métodos de cultivo solo logran detectar entre el 1 \% y el 5 \% de la comunidad microbiana total \cite{li2018deterioration}. Un estudio ilustrativo realizado en el campo de concentración de Auschwitz demostró que el cultivo detectó apenas 8 géneros bacterianos, mientras que la secuenciación NGS reveló 99 géneros en el mismo material \cite{gutarowska2015}. Adicionalmente, los medios de cultivo imponen concentraciones artificialmente elevadas de nutrientes, lo que favorece a microorganismos heterótrofos de rápido crecimiento que pueden no ser los colonizadores dominantes ni los principales agentes de deterioro in situ \cite{branysova2022, zanardini2019}.

Ahora bien, esta misma selectividad adquiere valor metodológico cuando el objetivo no es
inventariar la comunidad total, sino identificar a los organismos con capacidad real de
deterioro. A diferencia del ADN ambiental, que persiste en el sustrato tras la muerte
celular y no distingue entre células activas y muertas ~\cite{branysova2022}, el crecimiento
en cultivo constituye una prueba directa de viabilidad: todo microorganismo aislado está,
por definición, vivo. Por ello, la secuenciación metagenómica aplicada sobre la fracción
cultivada permite acoplar la certeza de viabilidad con la resolución genómica, estrategia
adoptada en la presente tesis~\cite{tsukayama2025informe}. La contrapartida, que se
mantiene explícita, es que dicha fracción representa únicamente el subconjunto cultivable
de la comunidad viable, por lo que sus conclusiones no se extienden a los linajes que no
crecen en los medios empleados.

\subsection{Tratamientos y limpieza con biocidas químicos}
\label{subsec:biocidas_quimicos}

Para el control del biodeterioro, la práctica conservativa tradicional recurrió históricamente a sustancias antimicrobianas de síntesis química: alcoholes, fenoles, aldehídos, compuestos de amonio cuaternario (como el cloruro de benzalconio), sales de estaño, derivados de isotiazolinonas y tratamientos de fumigación con óxido de etileno \cite{sterflinger2013, cappitelli2020}. Estas estrategias presentan ventajas a corto plazo en términos de reducción de biomasa visible, pero exhiben graves inconvenientes a mediano y largo plazo.

En primer lugar, la eficacia del biocida es frecuentemente transitoria y requiere reaplicaciones periódicas \cite{cappitelli2020}. En segundo lugar, generan toxicidad tanto para los operadores humanos como para el ecosistema circundante \cite{cappitelli2020}. Desde el punto de vista material, los biocidas oxidantes pueden inducir manchas por herrumbre y eflorescencias sobre la piedra, mientras que las soluciones de limpieza ácida o alcalina comprometen la integridad mineralógica del sustrato \cite{cappitelli2020}. Finalmente, y de mayor relevancia para la gestión del biodeterioro a largo plazo, los biocidas eliminan el microbioma natural sin remover la biomasa celular residual, la cual actúa como sustrato nutritivo que acelera la recolonización. Bajo la presión selectiva del biocida, las comunidades que recolonizan son frecuentemente resistentes al tratamiento y presentan perfiles metabólicos más agresivos, involucrando tanto a cepas bacterianas resistentes como a una mayor presencia de hongos altamente melanizados \cite{cappitelli2020, gorbushina2007}. Este efecto fue documentado de forma notable tras los tratamientos masivos realizados en las cuevas de Lascaux y en la tumba de San Pablo en Éfeso \cite{sterflinger2013}.

\section{Análisis comparativo de tecnologías actuales}
\label{sec:analisis_molecular_ngs}

La introducción de las tecnologías de secuenciación de próxima generación (NGS) como Illumina MiSeq, NovaSeq, PacBio y Oxford Nanopore, redefinió los estándares del análisis del microbioma del patrimonio pétro. En la práctica contemporánea, coexisten diversos enfoques que compiten en términos de rendimiento de datos, costos operativos y viabilidad de estandarización.

La secuenciación de amplicones de genes marcadores evolutivamente conservados, como las regiones hipervariables V3-V6 del gen 16S rRNA (para bacterias y arqueas) y las regiones ITS1/ITS2 (para hongos), permite la identificación filogenética de alta resolución sin necesidad de cultivo previo \cite{suchy2024, chimienti2016}.

La metagenómica de escopeta (shotgun metagenomics) secuencia la totalidad del material genético presente en la muestra, superando el sesgo de la amplificación específica por PCR y revelando no solo qué organismos están presentes, sino que permite determinar potencialmente qué rutas metabólicas y genes funcionales contiene la comunidad \cite{suchy2024, ding2020preahvihear}. Estudios previos demostraron que la secuenciación de amplicones permite detectar una diversidad microbiana significativamente superior a la obtenida mediante métodos clásicos y, en combinación con herramientas bioinformáticas predictivas, inferir el potencial funcional del microbioma en monumentos pétreos \cite{li2018deterioration, chimienti2016}.

Cabe precisar que la metagenómica \textit{shotgun} puede aplicarse tanto sobre ADN
extraído directamente del ambiente como sobre material previamente enriquecido por
cultivo; en el segundo caso, la comunidad se acota a su fracción viable a cambio de un
ADN de mayor calidad y trazabilidad, un compromiso relevante en muestras de baja biomasa
como las superficies pétreas~\cite{tsukayama2025informe,ding2020}.

En términos comparativos con los métodos tradicionales, la NGS supera significativamente al cultivo en cobertura taxonómica, sensibilidad y representatividad de la comunidad real \cite{sterflinger2013, suchy2024, gutarowska2015}. No obstante, su coste por muestra es considerablemente más elevado, requiere equipamiento especializado, genera grandes volúmenes de datos que demandan análisis bioinformáticos complejos, y presenta limitaciones propias: artefactos genéticos generados durante la amplificación por PCR (secuencias quiméricas), resolución taxonómica que frecuentemente no alcanza el nivel de especie, y dificultades de extracción de ADN en estructuras resistentes como las endosporas de Bacillus spp. \cite{branysova2022, li2018deterioration, skipper2022}.

En la frontera tecnológica se ubican la metatranscriptómica y la metabolómica ambiental, diseñadas para resolver el perfil funcional dinámico del microbioma in situ \cite{cappitelli2020, gutarowska2015}. La metatranscriptómica ofrece la medición directa de los transcritos de ARN activos al momento del muestreo, discriminando con precisión entre células muertas, en estado viable pero no cultivable (VBNC) o metabólicamente activas \cite{branysova2022}. Sin embargo, la secuenciación de ARN ambiental en el patrimonio pétreo se encuentra aún en una etapa inicial y carece de estandarización operativa; esto se debe en gran medida a la inestabilidad de la molécula y a que los costos de este análisis molecular siguen siendo demasiado altos en relación con los recursos habituales disponibles para los laboratorios de restauración \cite{suchy2024, sterflinger2013}. Por su parte, la metabolómica acoplada a espectrometría de masas (UPLC/HRMS) permite identificar los productos reales del catabolismo (como los ácidos orgánicos implicados en la disolución mineral), pero se enfrenta a la limitación de utilizar bases de datos de referencia ambiental incompletas que generan un volumen considerable de falsos positivos, situando los resultados en un nivel de identificación que requiere precaución y validación adicional \cite{gutarowska2015}.

\section{Insuficiencia tecnológica}
\label{sec:insuficiencia_tecnologica}

A pesar de los importantes avances moleculares, las soluciones actuales presentan vacíos estructurales al ser incapaces de confirmar si el potencial genético detectado se traduce en una actividad de deterioro real in situ, ya que el ADN perdura tras la muerte celular y disocia el inventario microbiano del daño activo \cite{ding2020preahvihear, cappitelli2020, branysova2022}. Además, existe una carencia crítica de metodologías estandarizadas y cuantitativas que permitan relacionar el perfil metabólico con el riesgo material \cite{cappitelli2020}, lo que impide diferenciar empíricamente si una biopelícula ejerce un rol de degradación o de bioprotección sobre la piedra \cite{cappitelli2020}. En este contexto, se ha subrayado la necesidad de transitar desde diagnósticos descriptivos y centrados en el material hacia marcos predictivos basados en la función microbiana \cite{gallego2025}; sin embargo, estas propuestas permanecen mayoritariamente en un plano conceptual, con énfasis en sustratos calcáreos y sin una implementación operativa que permita particionar el repertorio génico de los linajes colonizadores ni distinguir los rasgos de deterioro intrínsecos de aquellos adquiridos por adaptación al nicho. Por consiguiente, establecer una causalidad directa entre microorganismos específicos y el daño observado sigue siendo un reto profundo debido a las barreras económicas, logísticas y metodológicas para integrar datos multiómicos y petrofísicos \citep{branysova2022, gutarowska2015}, lo que evidencia la necesidad urgente de desarrollar plataformas diagnósticas integrales que vinculen la taxonomía con la actividad metabólica activa y el impacto real sobre el sustrato \cite{cappitelli2020}.